\section{Local Ramification Preliminaries}
This section reviews the ramification of discrete valuations, beginning with the general theory of \textit{discrete valuation rings} and their 
integral closures in finite separable extensions. 
We then specialize to complete rings in Galois extensions to describe the ramification filtration via lower and upper numbering, 
following the treatments in \stackstag{0EXQ} and \cite{serreLF}.

\begin{definition}[\stackstag{09E4}]\label{definition:weakly_unramified_and_tamely_ramified_DVR}
    We say that $A \to B$ or $A \subset B$ is an extension of discrete valuation rings if $A$ and $B$ are discrete valuation rings and $A \to B$ is injective and local. In particular, if $\pi _ A$ and $\pi _ B$ are uniformizers of $A$ and $B$, then $\pi _ A = u \pi _ B^ e$ for some $e \geq 1$ and unit $u$ of $B$. The integer $e$ does not depend on the choice of the uniformizers as it is also the unique integer $\geq 1$ such that
\[ \mathfrak m_ A B = \mathfrak m_ B^ e \]
The integer $e$ is called the \textit{ramification index} of $B$ over $A$. We say that $B$ is \textit{weakly unramified} over $A$ if $e = 1$. 
If the extension of residue fields $\kappa _ A = A/\mathfrak m_ A \subset \kappa _ B = B/\mathfrak m_ B$ is finite, 
then we set $f = [\kappa _ B : \kappa _ A]$ and we call it the \textit{residual degree} or residue degree of the extension $A \subset B$.

Note that we do not require the extension of fraction fields to be finite.
\end{definition}

Now let $A$ be a discrete valuation ring with fraction field $K$, let  $L/K$ be a finite separable field extension and 
let $B \subset L$ be the integral closure of $A$ in $L$.
% \[ \xymatrix{ B \ar[r] & L \\ A \ar[u] \ar[r] & K \ar[u] } \]
Then the ring extension $A \subset B$ is finite, hence $B$ is Noetherian. 
The dimension of $B$ is $1$, hence $B$ is a Dedekind domain.
Let $\mathfrak m_1, \ldots , \mathfrak m_ n$ be the maximal ideals of $B$ (i.e., the primes lying over $\mathfrak m_ A$). 
We obtain extensions of discrete valuation rings

\[ A \subset B_{\mathfrak m_ i} \]
and hence ramification indices $e_ i$ and residue degrees $f_ i$. We have

\[ [L : K] = \sum \nolimits _{i = 1, \ldots , n} e_ i f_ i \]
Note that if $A$ is henselian (e.g. $A$ is complete) then $n=1$. 

\begin{definition}[\stackstag{09E9}]\label{definition:tamely_ramified_dvrs}
Let $A$ be a discrete valuation ring with fraction field $K$. Let $L/K$ be a finite separable extension. With $B$ and $\mathfrak m_ i$, $i = 1, \ldots , n$ 
as above, we say the extension $L/K$ is
\begin{enumerate}
    \item unramified with respect to $A$ if $e_ i = 1$ and the extension $\kappa (\mathfrak m_ i)/\kappa _ A$ is separable for all $i$,
    \item tamely ramified with respect to $A$ if either the characteristic of $\kappa _ A$ is $0$ or the characteristic of $\kappa _ A$ is $p > 0$, the field extensions $\kappa (\mathfrak m_ i)/\kappa _ A$ are separable, and the ramification indices $e_ i$ are prime to $p$, and
    \item totally ramified with respect to $A$ if $n = 1$ and the residue field extension $\kappa (\mathfrak m_1)/\kappa _ A$ is trivial.
\end{enumerate}
If the discrete valuation ring $A$ is clear from context, then we sometimes say $L/K$ is unramified, totally ramified, or tamely ramified for short.
\end{definition}

% - %
Now and for the rest of the section assume $A$ is a \textit{complete discrete valuation ring} with uniformizer $\pi$ 
and residue field $\kappa$, which we assume to be perfect. When $A$ and $\kappa$ are of the same characteristic $p > 0$, 
then $A$ contains a coefficient field $k \cong \kappa$ and a well known structure theorem holds: $A = k[[\pi ]] \cong k[[t]]$.
Let $K$ be the fraction field of $A$, then $K=k((\pi))$.

For the remainder of this section, assume $A$ is a \textit{complete discrete valuation ring} with uniformizer $\pi$ and a perfect residue 
field $\kappa$. In the equal characteristic case, where $\text{char}(A) = \text{char}(\kappa) = p > 0$, the 
Cohen Structure Theorem implies that $A$ contains a coefficient field $k \cong \kappa$, such that $A \cong k[[\pi]] \cong k[[t]]$. 
Consequently, the fraction field is given by $K = k((\pi))$.
Let $L/K$ be a finite Galois extension with Galois group $G$. 
The integral closure $B$ of $A$ in $L$ is itself a complete discrete valuation ring. 
Since the residue field $\kappa$ is perfect, the extension of residue fields is separable, and there exists a uniformizer 
$\Pi \in B$ such that $B = A[\Pi]$.

The \textit{lower numbering ramification filtration} of $G$, denoted by $(G_i)_{i \geq -1}$, is defined as:
\[ 
G_i = \{ \sigma \in G \mid v_B(\sigma (x) - x) \geq i + 1 \text{ for all } x \in B \} 
\]
where $v_B$ is the valuation on $L$ associated with $B$. In this indexing, $G_{-1} = G$ and $G_0$ 
corresponds to the \textit{inertia group} of the extension $L/K$. Note that $L/K$ is unramified if and only if 
$G_0 = \{1\}$, and it is tamely ramified if and only if $G_1 = \{1\}$.

A standard result in the theory of local fields(complete local fields with a discrete valuation and perfect residue field) 
ensures that the condition in the definition of $G_i$ need only be checked for the uniformizer $\Pi$ of $B$. 
Specifically, if we define the function
\[ 
i^L_K(\sigma) = v_B(\sigma(\Pi) - \Pi) 
\]
for $\sigma \in G$, then $G_i = \{ \sigma \in G \mid i^L_K(\sigma) \geq i + 1 \}$. 
The subgroups $G_i$ are normal in $G$ and become trivial for sufficiently large $i$.

Consider a tower of fields $K \subset E \subset L$, and let $H = \text{Gal}(L/E) \subset G$. The lower numbering is compatible with subgroups:
\[ 
G_i \cap H = H_i \quad \text{for all } i \geq -1 
\]
which follows from the identity $i^L_E = i^L_K|_{H}$. 
However, the lower numbering does not behave as simply with respect to quotients. If $H$ is normal in $G$, the image of $G_i$ in $G/H$ is generally not $(G/H)_i$. 
Instead, we have $G_i H / H = (G/H)_j$, where the index $j$ is given by the formula:
\[ 
j = \frac{1}{e_{L/E}} \sum_{\tau \in H} \min(i^L_K(\tau), i+1) - 1 
\]


In the literature, one reindexes the ramification groups by defining the Herbrand function $\phi_{L/K} : [-1, \infty ) \to [-1, \infty )$:
\[ \phi_{L/K}(i) = \frac{1}{e_{L/K}}\sum_{\sigma \in G} \min(i^L_K(\sigma), i+1)-1 = \int _0^i \frac{1}{[G_0 : G_t]} dt \]

This function is continuous, strictly increasing, and piecewise linear, making it a bijection. 
It satisfies the composition law $\phi_{L/K} = \phi_{E/K} \circ \phi_{L/E}$ for a tower of extensions $K \subset E \subset L$. 
Using this function, we define the \textit{upper numbering ramification groups} as:
\[ 
G^i = G_{\phi_{L/K}^{-1}(i)} 
\]
The primary advantage of this reindexing is its compatibility with quotients: for any normal subgroup $H \subset G$, we have:
\[ 
G^i H / H = (G/H)^i 
\]
for all $i \geq -1$.


%% [REVIEW: never used -- the structure of G_i/G_{i+1} is not invoked anywhere (tameness is always derived via Kummer/Artin-Schreier); consider deleting or reducing to a one-sentence remark]
\begin{prop}[\Cite{serreLF}, IV Corollary 3]
    the quotients $G_i/G_{i+1}$, $i \geq 1$, are abelian groups, and are direct products of cyclic groups of order $p$. 
    The group $G_1$ is a $p$-group.
\end{prop}


%% [REVIEW: never used -- same as the proposition above; nothing downstream uses the semi-direct product structure of G_0]
\begin{prop}[\Cite{serreLF}, IV Corollary 4]
    The inertia group $G_0$ has the following property:
    \textit{It is the semi-direct product of a cyclic group of order prime to $p$ with a normal subgroup whose order is a power of $p$.}
\end{prop}

% We deduce the following lemma, that will be helpful later:
% \begin{lemma}
% Let $L/K$ be finite galois extesntion with galois group $G=\Z/p^r\Z$
% Assume there exists intermediate local field $K < M < L$ such that $M/K$ is unramified.
% Then $L/K$ is unramified.
% \end{lemma}




\subsection*{Kummer and Artin-Schreier Theories}
For cyclic extensions, Kummer and Artin-Schreier theories provide explicit ways to calculate ramification. 
These results show how the valuation of a defining element $a \in K$ determines the ramification index and the jumps in the upper numbering filtration.
Throughout this section, let $K$ be a discrete valuation field with a perfect residue field $\kappa$ of characteristic $p > 0$.

\begin{theorem}[Ramification in Kummer Extensions, \cite{koch1997algebraic}, Proposition 1.83]\label{theorem:kummer_ramification}
%Let $K$ be a local field of characteristic $\text{char}(K) =p$  
Assume $K$ contains the $n$-th roots of unity $\mu_n$. 
Let $L/K$ be the extension given by the equation $X^n = a$ for some 
$a \in K^\times$ and denote by $G$ its Galois group. Then we have:

\begin{enumerate}
    \item If $v_K(a) \in n\mathbb{Z}$ and the image of $a\pi^{-v_K(a)}$ in the residue field $\kappa$ is an $n$-th power, the extension $L/K$ is trivial.
    \item If $v_K(a) \in n\mathbb{Z}$ and the image of $a\pi^{-v_K(a)}$ in the residue field $\kappa$ is not an $n$-th power, the extension $L/K$ is cyclic and unramified.
    \item If $v_K(a) \notin n\mathbb{Z}$, the extension $L/K$ is cyclic and ramified. Specifically, if $\gcd(|v_K(a)|, n) = 1$, the extension is totally ramified of degree $n$.
     Otherwise it has ramification index $\frac{n}{|gcd(v_K(a)|, n)}$
\end{enumerate}
Conversly, Kummer theory ensures that every cyclic extension of degree $n$, prime to $p$ of a field that contains $n$-th roots of unity, is of the above form.
Moreover, in the above we can always take $a \in \Oo_K$
\end{theorem}
Note that in the case of total ramification the extesntion is tamely ramified.


\begin{theorem}[Ramification in Artin-Schreier Extensions, \cite{Thomas2005}]\label{theorem:artin_schreier_ramification}
Let $\wp(x) = x^p - x$ be the Artin-Schreier opeartor.
Let $L/K$ be the extension given by the equation $X^p - X = a$ for some $a \in K$ and denote by $G$ its Galois group. 
Then we have:
\begin{enumerate}
    \item If $v_K(a) > 0$ or if $v_K(a) = 0$ and $a \in \wp(K)$, the extension $L/K$ is trivial.
    \item If $v_K(a) = 0$ and if $a \notin \wp(K)$, the extension $L/K$ is cyclic of degree $p$ and unramified.
    \item If $v_K(a) = -m < 0$ with $m \in \mathbb{Z}_{>0}$ and if $m$ is prime to $p$, the extension $L/K$ is cyclic of degree $p$ again and totally ramified. Moreover,
     its ramification groups are given by:
     $$G = G^{(-1)} = \dots = G^{(m)} \quad \text{and} \quad G^{(m+1)} = 1.$$
\end{enumerate}
Conversely, Artin-Schreier theory ensures that every cyclic extension of degree $p$ takes this form.
Moreover, in the above and under the isomorphism $K \cong k((t))$, one can always take $a$ of the form 
$c t^{-m} + a_{-m + 1} t^{-m + 1} + \cdots a_{-1}t^{-1} + a_0$. If $k$ is algebraically closed then there is a 
change of variables such that $a=u^{-m}$.
\end{theorem}



%%%% %%%% %%%% %%%% %%%% %%%% %%%% %%%%
%          WITT VECTORS
%%%% %%%% %%%% %%%% %%%% %%%% %%%% %%%%
\providecommand{\bW}{\mathbb{W}}

\subsection{Artin--Schreier--Witt Theory}
\phantomsection\label{subsection:witt_vector_prerequisites}

Artin--Schreier theory, together with the ramification formula of
\Cref{theorem:artin_schreier_ramification}, gives complete control over the cyclic
extensions of degree $p$ of a local field of characteristic $p$. To treat the groups
$G = \Z/p^r\Z$ we need the analogous control over cyclic extensions of degree $p^r$,
which is provided by Witt's generalization: the field $K$ is replaced by the ring
$W_r(K)$ of $p$-typical Witt vectors of length $r$, the operator $\wp(u) = u^p - u$ by
its Witt-vector analogue $\wp = F - \mathrm{id}$, and the classical ramification-break
formula by the Schmid--Witt formula. The underlying Witt-ring algebra is collected in
\Cref{appendix:witt_vector_algebra}; this subsection retains the ASW classification,
the identities used in the proof, and the ramification statements. Our references for
the general theory are \cite[II, \S6]{serreLF} and \cite{Rabinoff2014Witt}.

Throughout, $p$ is a fixed prime, and $r \geq 1$ is an integer. All the Witt vectors
appearing in this thesis are $p$-typical and of finite length.


\paragraph{Working Witt-vector notation.}
For an \(\F_p\)-algebra \(A\), write \(W_r(A)=A^r\) for the ring of
\(p\)-typical Witt vectors of length \(r\), with Witt addition and subtraction
denoted by \(\oplus\) and \(\ominus\). We use the componentwise Frobenius \(F\),
the Verschiebung \(V\), and the truncation \(t\). The construction from Witt
polynomials, the ghost-map arguments, and the proofs of the identities among
these operators are given in \Cref{appendix:witt_vector_algebra}. In particular,
we use the triangular form of Witt addition
(\Cref{eq:witt_addition_shape}), the identities
\[
tW_r(A)=W_{r-1}(A),\qquad
\ker(t)=V^{r-1}A,\qquad
FV=VF=p,
\]
and the canonical identification \(W_r(\F_p)\cong\Z/p^r\Z\)
(\Cref{prop:witt_multiplication_by_p,example:witt_of_Fp}).

\begin{definition}\label{definition:asw_operator}
The \textbf{Artin--Schreier--Witt operator} on $W_r(A)$, $A$ an $\F_p$-algebra, is
\[
\wp \;:=\; F - \mathrm{id},
\qquad
\wp\,\vec u \;=\; F\vec u \ominus \vec u .
\]
For $r = 1$ this is the classical operator $\wp(u) = u^p - u$ of
\Cref{theorem:artin_schreier_ramification}.
\end{definition}

Since $F$ is a ring endomorphism, $\wp$ is an endomorphism of the additive group
$(W_r(A), \oplus)$:
$\wp(\vec u \oplus \vec v) = \wp\vec u \oplus \wp\vec v$. Moreover $\wp$ commutes with
$W_r(\sigma)$ for every ring homomorphism $\sigma$ (both $F$ and the ring operations
are natural), and with $V$ and $t$ by \Cref{prop:witt_multiplication_by_p}.

The following bookkeeping lemma is elementary but does the essential work of reducing
statements about $W_r$ to statements about $W_{r-1}$ and about $A$ itself: the
truncation $t$ and the inclusion $V^{r-1}$ of its kernel slice $W_r$ into lower-length
pieces, compatibly with $\wp$, and \emph{addition of an element of the kernel is
carry-free}.

\begin{lemma}[Witt bookkeeping]\label{lemma:witt_bookkeeping}
Let $A$ be an $\F_p$-algebra.
\begin{enumerate}[label=(\roman*)]
    \item $t \colon W_r(A) \to W_{r-1}(A)$ is a ring homomorphism commuting with $F$,
    hence with $\wp$; its kernel is $V^{r-1} A$.
    \item $V^{r-1} \colon A \to W_r(A)$ is additive, and
    $\wp \circ V^{r-1} = V^{r-1} \circ \wp$, where on the left $\wp$ is the operator
    on $W_r(A)$ and on the right the classical $\wp(c) = c^p - c$ on $A$.
    \item (carry-free top-level addition) For all $\vec a \in W_r(A)$ and $c \in A$,
    \[
    \vec a \oplus V^{r-1}(c) \;=\; (a_0, \dots, a_{r-2},\, a_{r-1} + c).
    \]
\end{enumerate}
\end{lemma}

\begin{proof}
(i) and (ii) restate parts of \Cref{prop:witt_multiplication_by_p} (for (ii): $V^{r-1}$
is additive and commutes with the componentwise $F$, hence with
$\wp = F - \mathrm{id}$).
(iii) By \Cref{lemma:ghost_principle} it suffices to verify the identity over
$\Z[a_0, \dots, a_{r-1}, c]$, where the ghost map is injective. Put
$\vec b := (0, \dots, 0, c)$ and $\vec s := \vec a \oplus \vec b$. For $n < r-1$ we
have $w_n(\vec b) = 0$, so $w_n(\vec s) = w_n(\vec a)$, and by induction on $n$
(solving for $s_n$, using that $p$ is a non-zero-divisor) $s_n = a_n$ for
$n \leq r-2$. For $n = r-1$, $w_{r-1}(\vec b) = p^{r-1} c$ gives
$p^{r-1} s_{r-1} = p^{r-1} a_{r-1} + p^{r-1} c$, i.e.\ $s_{r-1} = a_{r-1} + c$.
\end{proof}

\begin{rem*}
\Cref{lemma:witt_bookkeeping}(iii) fails for the addition of vectors supported in
lower positions: in general the carry polynomials $C_n$ of
\Cref{eq:witt_addition_shape} are non-zero, and Witt addition is \emph{not}
componentwise. In particular, operations performed componentwise on a Witt vector ---
such as discarding the regular parts of Laurent-series components --- do \emph{not}
in general respect Witt addition, and will be carefully avoided.
\end{rem*}

\subsubsection*{The Artin--Schreier--Witt isogeny and its torsors}

The componentwise ring operations make $\A^r_{\F_p}$ a ring scheme, the
\textbf{Witt vector group scheme} $\bW_r$: its underlying scheme is
$\A^r_{\F_p} = \Spec \F_p[T_0, \dots, T_{r-1}]$, with the group law given by the
addition polynomials $S_n$, so that $\bW_r(A) = W_r(A)$ functorially. The operator
$\wp = F - \mathrm{id}$ is an endomorphism of the group scheme $\bW_r$, and it is the
Witt-vector instance of a \emph{Lang isogeny}: by \Cref{example:witt_of_Fp} and the
lemma below, it sits in a short exact sequence of group schemes over $\F_p$ for the
\'etale topology
\begin{equation}\label{eq:asw_lang_sequence}
0 \longrightarrow \Z/p^r\Z \cong W_r(\F_p)
\longrightarrow \bW_r \xrightarrow{\;\wp\;} \bW_r \longrightarrow 0 .
\end{equation}

The corresponding covers admit the following completely explicit description, which
is the form in which we use them. For $\vec g \in W_r(A)$ we write
$A[\vec X]/(\wp \vec X \ominus \vec g)$, $\vec X = (X_0, \dots, X_{r-1})$, for the
quotient of $A[X_0, \dots, X_{r-1}]$ by the ideal generated by the $r$ components of
the Witt vector $\wp \vec X \ominus \vec g \in W_r(A[X_0, \dots, X_{r-1}])$.

\begin{lemma}[ASW covers are \'etale torsors]\label{lemma:asw_etale_torsor}
Let $A$ be an $\F_p$-algebra and $\vec g \in W_r(A)$. Let
\[
B \;:=\; A[\vec X]\big/\bigl(\wp \vec X \ominus \vec g\bigr).
\]
Then $B$ is finite free over $A$ with basis
$\bigl\{\prod_{n} X_n^{a_n} : 0 \le a_n \le p-1\bigr\}$, and \'etale over $A$. The
translation action of $W_r(\F_p) \cong \Z/p^r\Z$,
\[
m \colon \vec X \longmapsto \vec X \oplus \underline m ,
\]
makes $\Spec B$ a $\Z/p^r\Z$-torsor over $\Spec A$.
\end{lemma}

\begin{proof}
By \Cref{eq:witt_addition_shape} and the componentwise formula for $F$, the $n$-th
defining equation has the triangular form
\[
X_n^p - X_n - g_n + C_n'(X_0, \dots, X_{n-1};\, g_0, \dots, g_{n-1}) \;=\; 0
\]
for universal integral polynomials $C_n'$. Triangularity gives the free basis; the
Jacobian matrix of the system is lower triangular with diagonal entries
$p X_n^{p-1} - 1 = -1$, a unit, so $B$ is \'etale over $A$. Finally, $\Spec B$ is the
pullback along $\vec g \colon \Spec A \to \bW_r$ of the covering
$\wp \colon \bW_r \to \bW_r$; by the triangular form just established (over
$A = \F_p[T_0,\dots,T_{r-1}]$, $\vec g = (T_0,\dots,T_{r-1})$), $\wp$ is a surjective
\'etale homomorphism of group schemes, and its kernel is
$\{\vec u : F \vec u = \vec u\} = W_r(\F_p)$ since $F$ is componentwise. A surjective
\'etale group homomorphism is a torsor under its kernel, and torsors pull back to
torsors; hence $\Spec B$ is a $W_r(\F_p)$-torsor over $\Spec A$.
\end{proof}

\subsubsection*{Artin--Schreier--Witt theory for fields}

Over a field, the sequence \Cref{eq:asw_lang_sequence} computes all cyclic extensions
of $p$-power degree, exactly as classical Artin--Schreier theory does for degree $p$.

\begin{theorem}[Artin--Schreier--Witt theory; {\cite[\S3--4]{Thomas2005}}]
\label{theorem:asw_theory}
Let $M \supseteq \F_p$ be a field, $M^{s}$ a separable closure, and
$\Gamma = \Gal(M^{s}/M)$. Then $\wp$ is surjective on $W_r(M^{s})$ with kernel
$W_r(\F_p)$, and there is a canonical isomorphism of abelian groups
\[
\delta \colon\; W_r(M)\big/\wp\, W_r(M) \;\xrightarrow{\;\cong\;}\;
\Hom_{\mathrm{cont}}\bigl(\Gamma,\, \Z/p^r\Z\bigr) \;=\; H^1\bigl(M, \Z/p^r\Z\bigr),
\qquad
[\vec g] \longmapsto \bigl(\sigma \mapsto \sigma(\vec u) \ominus \vec u\bigr),
\]
where $\vec u \in W_r(M^{s})$ is any solution of $\wp \vec u = \vec g$. Under the
correspondence between characters and torsors
(\Cref{cor:abelian_equivilance_fundamental_torsors}), the class $[\vec g]$ corresponds
to the $\Z/p^r\Z$-torsor
\[
\Spec\, M[\vec X]\big/\bigl(\wp \vec X \ominus \vec g\bigr)
\]
of \Cref{lemma:asw_etale_torsor}. Moreover, for a class of order $p^s$ in
$W_r(M)/\wp W_r(M)$, the fixed field $M(\vec u) := M(u_0, \dots, u_{r-1})$ of
$\ker \delta[\vec g]$ is a cyclic extension of $M$ of degree $p^s$; the torsor is
connected if and only if $[\vec g]$ has order $p^r$, in which case
$\Spec M[\vec X]/(\wp \vec X \ominus \vec g) \cong \Spec M(\vec u)$ is (the spectrum
of) a cyclic field extension of degree $p^r$. Every cyclic extension of $M$ of degree
dividing $p^r$ arises this way.
\end{theorem}

\begin{proof}[Sketch of proof]
Surjectivity of $\wp$ on $W_r(M^{s})$ and the computation of its kernel follow from
the triangular equations in the proof of \Cref{lemma:asw_etale_torsor}: each component
of a preimage is obtained by solving a separable equation $X_n^p - X_n = (\dots)$, and
$F\vec u = \vec u$ holds if and only if $u_n^p = u_n$ for all $n$, i.e.\
$\vec u \in W_r(\F_p)$. Given $\vec g$ and a solution $\wp\vec u = \vec g$, for
$\sigma \in \Gamma$ we have
$\wp(\sigma \vec u \ominus \vec u) = \sigma \vec g \ominus \vec g = 0$ (using
\Cref{eq:witt_invariants_componentwise} and that $W_r(\sigma)$ commutes with $\wp$),
so $\sigma \vec u \ominus \vec u \in W_r(\F_p)$; one checks it is a continuous
homomorphism in $\sigma$ (the kernel is fixed pointwise by $\Gamma$), independent of
the choices modulo nothing and of the representative $\vec g$ modulo $\wp W_r(M)$.
Injectivity: if $\delta[\vec g] = 0$ then $\vec u$ is $\Gamma$-invariant, hence lies
in $W_r(M)$ by \Cref{eq:witt_invariants_componentwise}, so
$\vec g = \wp\vec u \in \wp W_r(M)$. Surjectivity of $\delta$ is the additive
Hilbert~90: $H^1(\Gamma, W_r(M^{s})) = 0$, by induction on $r$ along the exact
sequence $0 \to M^{s} \xrightarrow{V^{r-1}} W_r(M^{s}) \xrightarrow{t} W_{r-1}(M^{s})
\to 0$ of $\Gamma$-modules (\Cref{lemma:witt_bookkeeping}(i)) and the classical case
$H^1(\Gamma, M^{s}) = 0$. The remaining statements follow from
\Cref{theorem:equivalence_fundamental_covers} applied to the character
$\delta[\vec g]$, whose image has the same order as the class $[\vec g]$. For details
see \cite[\S3--4]{Thomas2005} or \cite{Rabinoff2014Witt}.
\end{proof}

The order of a class is partially visible in its $0$-th component:

\begin{cor}\label{cor:asw_class_order_and_first_component}
Let $M \supseteq \F_p$ be a field and $\vec g \in W_r(M)$ with $g_0 = 0$. Then
$p^{r-1}[\vec g] = 0$ in $W_r(M)/\wp W_r(M)$; equivalently, the cyclic extension
attached to $[\vec g]$ has degree at most $p^{r-1}$. Consequently, if $[\vec g]$ has
order $p^r$ (e.g.\ if $\vec g$ defines a connected torsor), then \emph{every}
representative of the class has non-zero $0$-th component.
\end{cor}

\begin{proof}
By \Cref{eq:witt_p_to_r_minus_1}, $p^{r-1} \vec g = (0, \dots, 0, g_0^{p^{r-1}}) = 0$
already in $W_r(M)$.
\end{proof}

\begin{rem*}\label{rem:asw_local_totally_ramified}
When $M = k((x))$ with $k$ algebraically closed, every finite separable extension
$L/M$ is totally ramified: the residue field extension is trivial (as $k$ is
algebraically closed) and $\Oo_L$ is a complete discrete valuation ring, so
$[L : M] = ef = e$. Thus a class of order $p^r$ in $W_r(k((x)))/\wp$ corresponds to a
cyclic, \emph{totally ramified} extension of degree $p^r$.
\end{rem*}



The isobaricity property of the Witt addition polynomials used in the later
degree estimate is recorded, with its proof, in
\Cref{lemma:witt_addition_isobaric}.

\subsubsection*{Ramification breaks: the Schmid--Witt formula}

Finally we state the generalization of the ramification formula of
\Cref{theorem:artin_schreier_ramification} to Artin--Schreier--Witt extensions. As in
the classical case, the formula reads off the largest upper-numbering ramification
break from the pole orders of a defining Witt vector --- provided the representative
is normalized so that no interference between the levels can occur. The relevant
normalization is the following.

\begin{definition}\label{definition:reduced_witt_vector}
Let $K = k((x))$. For $g \in K$ put $m(g) := \max\bigl(0, -v_x(g)\bigr)$ (the pole
order of $g$, or $0$ if $g$ is regular). A vector
$\vec g = (g_0, \dots, g_{r-1}) \in W_r(K)$ is \textbf{reduced} if for every $j$,
either $m(g_j) = 0$ or $p \nmid m(g_j)$.
\end{definition}

\begin{theorem}[Schmid--Witt break formula]\label{theorem:schmid_witt_break_formula}
Let $K = k((x))$ with $k$ perfect of characteristic $p$, and let $L/K$ be a cyclic
totally ramified extension of degree $p^r$ whose Artin--Schreier--Witt class
(\Cref{theorem:asw_theory}) is represented by a \emph{reduced} vector
$\vec g \in W_r(K)$, with $m_j := m(g_j)$. Then the largest upper-numbering
ramification break of $L/K$ equals
\[
u \;=\; \max_{0 \le j \le r-1} \; p^{\,r-1-j}\, m_j .
\]
\end{theorem}

This is due, for finite residue fields, to Kanesaka--Sekiguchi
\cite{KanesakaSekiguchi1979} and Thomas \cite{Thomas2005} (building on Schmid
\cite{Schmid1937} for $r = 1$), and for arbitrary perfect residue fields to
Elder--Keating \cite{ElderKeating2025}. We use it as stated, in one direction only:
\emph{evaluating the formula on one reduced representative computes $u$}; no
minimization over representatives is needed.

\begin{rem*}
Reducedness is exactly what makes the formula exact. If two levels $j < j'$ with
$m_j, m_{j'} \neq 0$ satisfied $p^{r-1-j} m_j = p^{r-1-j'} m_{j'}$, then
$m_{j'} = p^{\,j'-j} m_j$ would be divisible by $p$, contradicting
\Cref{definition:reduced_witt_vector}; hence the maximum is attained at a unique
level, and no cancellation of breaks between levels can occur.
\end{rem*}

\begin{rem*}
For $r = 1$, a reduced vector is a single element $g \in K$ whose pole order $m_0$ is
zero or prime to $p$, and \Cref{theorem:schmid_witt_break_formula} recovers
\Cref{theorem:artin_schreier_ramification}: the largest (indeed the only) upper break
is $m_0$. The existence of reduced representatives --- and, over $K = k((x))$, of
reduced representatives whose components are polynomials in $x^{-1}$ without constant
term --- is not part of the general theory quoted here; it will be proved where it is
used, by induction on $r$ through the truncation maps of
\Cref{lemma:witt_bookkeeping}.
\end{rem*}


\section{A Base-Change Lemma for Galois Extensions}

Finally, we record a standard base-change lemma for Galois field extensions, used when
comparing a Galois extension with the function field of a base change.

\begin{lemma}\label{lem:galois-base-change}
Let $E/F$ be a finite Galois extension with group $H$, and let $K/F$ be an arbitrary
field extension. Fix a common overfield $\Omega \supseteq E,K$ (for instance an algebraic
closure of $K$ together with an $F$-embedding $E \hookrightarrow \Omega$), and form inside
it the intersection $M := E \cap K$, the compositum $E\cdot K$, and the subgroup
$N := \Gal(E/M) \leq H$. Then
\begin{enumerate}
  \item $E\cdot K / K$ is Galois, and restriction is an isomorphism
        $\Gal(E\cdot K/K) \xrightarrow{\ \sim\ } \Gal(E/M) = N$; and
  \item as $K$-algebras,
        $\displaystyle E \otimes_F K \;\cong\; \prod_{i=1}^{[M:F]} E\cdot K$.
\end{enumerate}
In particular $E\otimes_F K$ is a field if and only if $M = F$, i.e.\ if and only if $E$ and
$K$ are linearly disjoint over $F$.
\end{lemma}

\begin{proof}
Write $E = F[\alpha]$ and let $g \in F[x]$ be the minimal polynomial of $\alpha$, so
$\deg g = |H|$ and $g$ splits completely in $E$ because $E/F$ is Galois. Part (1) is the
classical translation theorem: $E\cdot K/K$ is Galois (it is generated over $K$ by the
roots of the separable $g$, which all lie in $E$), and $\sigma \mapsto \sigma|_E$ embeds
$\Gal(E\cdot K/K)$ into $H$ with image the subgroup fixing $E\cap K = M$ pointwise, namely
$\Gal(E/M) = N$. For (2), factor $g = g_1 \cdots g_m$ into monic irreducibles over $K$.
Adjoining any single root of $g$ to $K$ already produces all of $E\cdot K$ (since
$E = F(\alpha)$ contains every conjugate of $\alpha$), so each $g_i$ has degree
$[E\cdot K : K] = [E:M] = |N|$ by (1); hence $m = |H|/|N| = [M:F]$. Therefore
\[
E \otimes_F K \;\cong\; K[x]/(g) \;\cong\; \prod_{i=1}^{m} K[x]/(g_i)
\;\cong\; \prod_{i=1}^{[M:F]} E\cdot K ,
\]
the last step because $K[x]/(g_i) \cong K(\alpha) = E\cdot K$ for every $i$. This product is
a single field exactly when $m = 1$, i.e.\ when $[M:F] = 1$, i.e.\ when $M = F$.
\end{proof}
